%!TeX root=../persuasiontop.tex
\chapter[Chapter \thechapter]{}

\begin{figure}[t!]
\centering
\includegraphics[width=\linewidth]{nouveau19}
\end{figure}

\lettrine[lraise=0.3]{W}{hile} Admiral Croft was taking this walk with Anne, and expressing his wish of getting Captain Wentworth to Bath, Captain Wentworth was already on his way thither. Before Mrs Croft had written, he was arrived, and the very next time Anne walked out, she saw him.

Mr Elliot was attending his two cousins and Mrs Clay. They were in Milsom Street. It began to rain, not much, but enough to make shelter desirable for women, and quite enough to make it very desirable for Miss Elliot to have the advantage of being conveyed home in Lady Dalrymple's carriage, which was seen waiting at a little distance; she, Anne, and Mrs Clay, therefore, turned into Molland's, while Mr Elliot stepped to Lady Dalrymple, to request her assistance. He soon joined them again, successful, of course; Lady Dalrymple would be most happy to take them home, and would call for them in a few minutes.

Her ladyship's carriage was a barouche, and did not hold more than four with any comfort. Miss Carteret was with her mother; consequently it was not reasonable to expect accommodation for all the three Camden Place ladies. There could be no doubt as to Miss Elliot. Whoever suffered inconvenience, she must suffer none, but it occupied a little time to settle the point of civility between the other two. The rain was a mere trifle, and Anne was most sincere in preferring a walk with Mr Elliot. But the rain was also a mere trifle to Mrs Clay; she would hardly allow it even to drop at all, and her boots were so thick! much thicker than Miss Anne's; and, in short, her civility rendered her quite as anxious to be left to walk with Mr Elliot as Anne could be, and it was discussed between them with a generosity so polite and so determined, that the others were obliged to settle it for them; Miss Elliot maintaining that Mrs Clay had a little cold already, and Mr Elliot deciding on appeal, that his cousin Anne's boots were rather the thickest.

It was fixed accordingly, that Mrs Clay should be of the party in the carriage; and they had just reached this point, when Anne, as she sat near the window, descried, most decidedly and distinctly, Captain Wentworth walking down the street.

Her start was perceptible only to herself; but she instantly felt that she was the greatest simpleton in the world, the most unaccountable and absurd! For a few minutes she saw nothing before her; it was all confusion. She was lost, and when she had scolded back her senses, she found the others still waiting for the carriage, and Mr Elliot (always obliging) just setting off for Union Street on a commission of Mrs Clay's.

She now felt a great inclination to go to the outer door; she wanted to see if it rained. Why was she to suspect herself of another motive? Captain Wentworth must be out of sight. She left her seat, she would go; one half of her should not be always so much wiser than the other half, or always suspecting the other of being worse than it was. She would see if it rained. She was sent back, however, in a moment by the entrance of Captain Wentworth himself, among a party of gentlemen and ladies, evidently his acquaintance, and whom he must have joined a little below Milsom Street. He was more obviously struck and confused by the sight of her than she had ever observed before; he looked quite red. For the first time, since their renewed acquaintance, she felt that she was betraying the least sensibility of the two. She had the advantage of him in the preparation of the last few moments. All the overpowering, blinding, bewildering, first effects of strong surprise were over with her. Still, however, she had enough to feel! It was agitation, pain, pleasure, a something between delight and misery.

He spoke to her, and then turned away. The character of his manner was embarrassment. She could not have called it either cold or friendly, or anything so certainly as embarrassed.

After a short interval, however, he came towards her, and spoke again. Mutual enquiries on common subjects passed: neither of them, probably, much the wiser for what they heard, and Anne continuing fully sensible of his being less at ease than formerly. They had by dint of being so very much together, got to speak to each other with a considerable portion of apparent indifference and calmness; but he could not do it now. Time had changed him, or Louisa had changed him. There was consciousness of some sort or other. He looked very well, not as if he had been suffering in health or spirits, and he talked of Uppercross, of the Musgroves, nay, even of Louisa, and had even a momentary look of his own arch significance as he named her; but yet it was Captain Wentworth not comfortable, not easy, not able to feign that he was.

It did not surprise, but it grieved Anne to observe that Elizabeth would not know him. She saw that he saw Elizabeth, that Elizabeth saw him, that there was complete internal recognition on each side; she was convinced that he was ready to be acknowledged as an acquaintance, expecting it, and she had the pain of seeing her sister turn away with unalterable coldness.

Lady Dalrymple's carriage, for which Miss Elliot was growing very impatient, now drew up; the servant came in to announce it. It was beginning to rain again, and altogether there was a delay, and a bustle, and a talking, which must make all the little crowd in the shop understand that Lady Dalrymple was calling to convey Miss Elliot. At last Miss Elliot and her friend, unattended but by the servant, (for there was no cousin returned), were walking off; and Captain Wentworth, watching them, turned again to Anne, and by manner, rather than words, was offering his services to her.

<I am much obliged to you,> was her answer, <but I am not going with them. The carriage would not accommodate so many. I walk: I prefer walking.>

<But it rains.>

<Oh! very little. Nothing that I regard.>

After a moment's pause he said: <Though I came only yesterday, I have equipped myself properly for Bath already, you see,> (pointing to a new umbrella); <I wish you would make use of it, if you are determined to walk; though I think it would be more prudent to let me get you a chair.>

She was very much obliged to him, but declined it all, repeating her conviction, that the rain would come to nothing at present, and adding, <I am only waiting for Mr Elliot. He will be here in a moment, I am sure.>

She had hardly spoken the words when Mr Elliot walked in. Captain Wentworth recollected him perfectly. There was no difference between him and the man who had stood on the steps at Lyme, admiring Anne as she passed, except in the air and look and manner of the privileged relation and friend. He came in with eagerness, appeared to see and think only of her, apologised for his stay, was grieved to have kept her waiting, and anxious to get her away without further loss of time and before the rain increased; and in another moment they walked off together, her arm under his, a gentle and embarrassed glance, and a <Good morning to you!> being all that she had time for, as she passed away.

As soon as they were out of sight, the ladies of Captain Wentworth's party began talking of them.

<Mr Elliot does not dislike his cousin, I fancy?>

<Oh! no, that is clear enough. One can guess what will happen there. He is always with them; half lives in the family, I believe. What a very good-looking man!>

<Yes, and Miss Atkinson, who dined with him once at the Wallises, says he is the most agreeable man she ever was in company with.>

<She is pretty, I think; Anne Elliot; very pretty, when one comes to look at her. It is not the fashion to say so, but I confess I admire her more than her sister.>

<Oh! so do I\@.>

<And so do I\@. No comparison. But the men are all wild after Miss Elliot. Anne is too delicate for them.>

Anne would have been particularly obliged to her cousin, if he would have walked by her side all the way to Camden Place, without saying a word. She had never found it so difficult to listen to him, though nothing could exceed his solicitude and care, and though his subjects were principally such as were wont to be always interesting: praise, warm, just, and discriminating, of Lady Russell, and insinuations highly rational against Mrs Clay. But just now she could think only of Captain Wentworth. She could not understand his present feelings, whether he were really suffering much from disappointment or not; and till that point were settled, she could not be quite herself.

She hoped to be wise and reasonable in time; but alas! alas! she must confess to herself that she was not wise yet.

Another circumstance very essential for her to know, was how long he meant to be in Bath; he had not mentioned it, or she could not recollect it. He might be only passing through. But it was more probable that he should be come to stay. In that case, so liable as every body was to meet every body in Bath, Lady Russell would in all likelihood see him somewhere. Would she recollect him? How would it all be?

She had already been obliged to tell Lady Russell that Louisa Musgrove was to marry Captain Benwick. It had cost her something to encounter Lady Russell's surprise; and now, if she were by any chance to be thrown into company with Captain Wentworth, her imperfect knowledge of the matter might add another shade of prejudice against him.

The following morning Anne was out with her friend, and for the first hour, in an incessant and fearful sort of watch for him in vain; but at last, in returning down Pulteney Street, she distinguished him on the right hand pavement at such a distance as to have him in view the greater part of the street. There were many other men about him, many groups walking the same way, but there was no mistaking him. She looked instinctively at Lady Russell; but not from any mad idea of her recognising him so soon as she did herself. No, it was not to be supposed that Lady Russell would perceive him till they were nearly opposite. She looked at her however, from time to time, anxiously; and when the moment approached which must point him out, though not daring to look again (for her own countenance she knew was unfit to be seen), she was yet perfectly conscious of Lady Russell's eyes being turned exactly in the direction for him—of her being, in short, intently observing him. She could thoroughly comprehend the sort of fascination he must possess over Lady Russell's mind, the difficulty it must be for her to withdraw her eyes, the astonishment she must be feeling that eight or nine years should have passed over him, and in foreign climes and in active service too, without robbing him of one personal grace!

At last, Lady Russell drew back her head. <Now, how would she speak of him?>

<You will wonder,> said she, <what has been fixing my eye so long; but I was looking after some window-curtains, which Lady Alicia and Mrs Frankland were telling me of last night. They described the drawing-room window-curtains of one of the houses on this side of the way, and this part of the street, as being the handsomest and best hung of any in Bath, but could not recollect the exact number, and I have been trying to find out which it could be; but I confess I can see no curtains hereabouts that answer their description.>

Anne sighed and blushed and smiled, in pity and disdain, either at her friend or herself. The part which provoked her most, was that in all this waste of foresight and caution, she should have lost the right moment for seeing whether he saw them.

A day or two passed without producing anything. The theatre or the rooms, where he was most likely to be, were not fashionable enough for the Elliots, whose evening amusements were solely in the elegant stupidity of private parties, in which they were getting more and more engaged; and Anne, wearied of such a state of stagnation, sick of knowing nothing, and fancying herself stronger because her strength was not tried, was quite impatient for the concert evening. It was a concert for the benefit of a person patronised by Lady Dalrymple. Of course they must attend. It was really expected to be a good one, and Captain Wentworth was very fond of music. If she could only have a few minutes conversation with him again, she fancied she should be satisfied; and as to the power of addressing him, she felt all over courage if the opportunity occurred. Elizabeth had turned from him, Lady Russell overlooked him; her nerves were strengthened by these circumstances; she felt that she owed him attention.

She had once partly promised Mrs Smith to spend the evening with her; but in a short hurried call she excused herself and put it off, with the more decided promise of a longer visit on the morrow. Mrs Smith gave a most good-humoured acquiescence.

<By all means,> said she; <only tell me all about it, when you do come. Who is your party?>

Anne named them all. Mrs Smith made no reply; but when she was leaving her said, and with an expression half serious, half arch, <Well, I heartily wish your concert may answer; and do not fail me to-morrow if you can come; for I begin to have a foreboding that I may not have many more visits from you.>

Anne was startled and confused; but after standing in a moment's suspense, was obliged, and not sorry to be obliged, to hurry away.